% ============================================================
% Preambulo moderno para Tarea Práctica 1
% Basado en la plantilla original del autor.
% Se conservan las interfaces: \materia, \TituloTarea, \Tony y \titulo.
% ============================================================

% Asegura que los .sty locales se encuentren dentro de la carpeta plantilla/
\makeatletter
%\def\input@path{{plantilla/}}
\makeatother

\usepackage{fontspec}      % Si no lo tienes, añádelo antes de babel
\usepackage[spanish,es-tabla,activeacute]{babel}
\selectlanguage{spanish}
\spanishdecimal{.}

\usepackage{varios}
\usepackage{needspace}  % evita que tablas cortas se partan entre páginas
\usepackage{tikz}       % portada ejecutiva (bloques y KPIs)
\graphicspath{{plantilla/}}
% ------------------------------------------------------------------
% Paleta Línea 12 (Línea Dorada) — sobrescribe los colores UNAM.
% Todo lo que ya usa estos nombres (portada, secciones, cajas,
% hipervínculos) hereda el dorado automáticamente.
% ------------------------------------------------------------------
\definecolor{UNAMBlue}{HTML}{006199}        % dorado principal
\definecolor{UNAMAccent}{HTML}{8A6A16}      % dorado oscuro (contraste)
\definecolor{UNAMLightBlue}{HTML}{F3E7C6}   % fondo claro dorado
\definecolor{UNAMLightAccent}{HTML}{EAD6A0} % fondo claro dorado oscuro
\definecolor{UNAMMuted}{HTML}{5B5B5B}       % gris neutro, sin cambio
\definecolor{UNAMText}{HTML}{262626}        % texto principal
\definecolor{UNAMLightGray}{HTML}{FBF9F4}   % blanco hueso

% Matemáticas modernas para XeLaTeX
\usepackage{unicode-math}
\setmathfont{TeX Gyre Termes Math}

\renewcommand{\labelitemi}{$\bullet$}
\renewcommand{\labelitemii}{$\circ$}

% ============================================================
% Identidad del documento: conservar estos comandos
% ============================================================
\newcommand{\materia}{Modelos de Supervivencia y Series de Tiempo}
\newcommand{\TituloTarea}{Retención de clientes en Divvy}
%\newcommand{\Tony}{\textsc{Anthonie Byron Ornelas Clement}}

% Autores
\newcommand{\Autores}{
    \textsc{Aguirre Garcia Abraham Isaias}\\
    \textsc{Castillo C\'ordoba Alfa Cecilia}\\
    \textsc{Larios Huitron Michelle}\\
    \textsc{Ornelas Clement Anthonie Byron}
}

% ============================================================
% Portada (opción 2 — reporte de decisión con KPIs)
% NOTA: las cifras de los KPIs provienen del análisis del Rmd;
% si la base o el modelo cambian, sincronizarlas a mano aquí.
% ============================================================
\definecolor{PortNavy}{HTML}{264653}
\definecolor{PortTeal}{HTML}{2A9D8F}
\definecolor{PortFuga}{HTML}{E63946}
\definecolor{PortAcento}{HTML}{F77F00}
\definecolor{PortMuted}{HTML}{65727A}

\newcommand{\titulo}{%
  \begin{tikzpicture}[remember picture, overlay]
    \fill[PortNavy] (current page.north west) rectangle ([yshift=-8.5cm]current page.north east);
    \node[anchor=north west, text width=17.8cm] at ([xshift=1.6cm,yshift=-1.5cm]current page.north west) {
      {\sffamily\fontsize{13}{16}\selectfont\color{white!75!PortNavy}INFORME EJECUTIVO DE RETENCIÓN}\\[0.45cm]
      {\sffamily\fontsize{36}{42}\selectfont\bfseries\color{white} \TituloTarea\\[0.3cm]
       \fontsize{15}{20}\selectfont\mdseries ¿Qué modificación al producto mejora la retención de clientes?}
    };
    \node[anchor=north west, text width=17.8cm] at ([xshift=1.6cm,yshift=-9.6cm]current page.north west) {
      {\sffamily\small\color{PortMuted}LO QUE ENCONTRAMOS EN 3 NÚMEROS}\\[0.35cm]
      \begin{tabular}{c}
        {\sffamily\fontsize{26}{30}\selectfont\bfseries\color{PortNavy}1,000}\\[-0.1cm]{\sffamily\footnotesize\color{PortMuted}clientes analizados}
      \end{tabular}\hspace{1.6cm}
      \begin{tabular}{c}
        {\sffamily\fontsize{26}{30}\selectfont\bfseries\color{PortFuga}52.6\%}\\[-0.1cm]{\sffamily\footnotesize\color{PortMuted}de fuga en el año}
      \end{tabular}\hspace{1.6cm}
      \begin{tabular}{c}
        {\sffamily\fontsize{26}{30}\selectfont\bfseries\color{PortTeal}269 días}\\[-0.1cm]{\sffamily\footnotesize\color{PortMuted}mediana de fuga}
      \end{tabular}\\[0.9cm]
      {\sffamily\color{PortTeal}\rule{17.8cm}{1.6pt}}\\[0.35cm]
      {\sffamily\small\color{PortNavy}\bfseries Decisión que este informe sostiene:}\\[-3mm]
      {\sffamily\small\color{PortNavy}ciclo de re-engagement universal antes del día 180 y re-instrumentación del registro.}\\[1.0cm]
      {\sffamily\color{PortNavy}\Autores}
    };
    \node[anchor=south east] at ([xshift=-1.4cm,yshift=1.2cm]current page.south east) {\includegraphics[height=2.4cm]{escudo0.png}};
  \end{tikzpicture}%
}

% ============================================================
% Ajustes visuales generales
% ============================================================
\setlength{\parindent}{0pt}
\setlength{\parskip}{0.55em}

% Hyperref lo administra Pandoc/R Markdown. Solo configuramos al inicio.
\makeatletter
\AtBeginDocument{%
  \@ifpackageloaded{hyperref}{%
    \hypersetup{
      colorlinks=true,
      linkcolor=UNAMBlue,
      urlcolor=UNAMBlue,
      citecolor=UNAMBlue
    }
  }{}
}
\makeatother

% Títulos
\usepackage{titlesec}
\titleformat{\section}
  {\sffamily\Large\bfseries\color{UNAMBlue}}
  {\thesection}{0.65em}{}
  [\vspace{0.18em}\color{UNAMAccent}\titlerule]

\titleformat{\subsection}
  {\sffamily\large\bfseries\color{UNAMText}}
  {\thesubsection}{0.6em}{}

\titleformat{\subsubsection}
  {\sffamily\normalsize\bfseries\color{UNAMBlue}}
  {\thesubsubsection}{0.55em}{}

% Figuras
\captionsetup{
  font=small,
  labelfont={bf,color=UNAMBlue},
  textfont={color=UNAMMuted},
  justification=centering,
  skip=7pt
}

% ============================================================
% Entornos de comunicación para el Rmd
% ============================================================
\newtcolorbox{queobservar}{%
  enhanced,
  breakable,
  colback=UNAMLightBlue,
  colframe=UNAMBlue,
  boxrule=0pt,
  leftrule=3pt,
  arc=2pt,
  left=8pt,right=8pt,top=7pt,bottom=7pt,
  before skip=8pt,after skip=10pt,
  fonttitle=\sffamily\bfseries,
  title=¿Qué debemos observar?%
}

\newtcolorbox{notatecnica}{%
  enhanced,
  breakable,
  colback=UNAMLightGray,
  colframe=UNAMMuted,
  boxrule=0pt,
  leftrule=3pt,
  arc=2pt,
  left=8pt,right=8pt,top=7pt,bottom=7pt,
  before skip=8pt,after skip=10pt,
  fonttitle=\sffamily\bfseries,
  title=Nota técnica%
}

\newtcolorbox{hallazgo}{%
  enhanced,
  breakable,
  colback=UNAMLightAccent,
  colframe=UNAMAccent,
  boxrule=0pt,
  leftrule=4pt,
  arc=2pt,
  left=8pt,right=8pt,top=8pt,bottom=8pt,
  before skip=9pt,after skip=11pt,
  fonttitle=\sffamily\bfseries,
  title=Hallazgo principal%
}


\newtcolorbox{citadestacada}{%
  enhanced,
  breakable,
  colback=white,
  colframe=UNAMBlue,
  boxrule=0pt,
  leftrule=2.4pt,
  arc=0pt,
  left=12pt,right=10pt,top=6pt,bottom=6pt,
  before skip=8pt,after skip=10pt,
  fontupper=\itshape\color{UNAMText}
}